\section*{Appendix F — Meaning Ledger Formalism}
\addcontentsline{toc}{section}{Appendix F — Meaning Ledger Formalism}

\noindent This appendix formalizes the Meaning Ledger, the substrate’s global
justification graph. The ledger records all lawful semantic, runtime, and
interaction transitions, enforces invariant preservation, detects drift,
regulates load, and ensures curvature continuity across the meaning manifold
$\mathcal{M}$. The ledger is the substrate’s global legality substrate: no
transition may enter the cognitive state unless it is justified, invariant-
preserving, and manifold-legal.

\medskip

\noindent The v0.6 Meaning Ledger consists of the following components:

\begin{itemize}
    \item F.1 Ledger Graph Structure (Justification Graph Geometry)
    \item F.2 Invariant Preservation (Meaning, Identity, Geometry, Load)
    \item F.3 Drift Detection Metrics (Curvature, Identity, Geometry, Load)
    \item F.4 Ledger-Closed Reasoning (Legality-Preserved Cognitive Motion)
    \item F.5 Cosmological Ledger Geometry (Multi-Manifold Ledger Coupling)
    \item F.6 Collapse & Capture Signatures (Ledger-Based Failure Detection)
    \item F.7 Restoration Signatures (Ledger-Based Recovery Geometry)
    \item F.8 Ledger Closure & Operator Algebra (Extended Closure Conditions)
\end{itemize}

% ------------------------------------------------------------
% F.1 LEDGER GRAPH STRUCTURE
% ------------------------------------------------------------

\subsection*{F.1 Ledger Graph Structure (Justification Graph Geometry)}

\noindent The Meaning Ledger is the substrate’s global justification graph. It
records all lawful semantic, runtime, and interaction transitions across the
meaning manifold $\mathcal{M}$, ensuring that every cognitive step is invariant-
preserving, curvature-stable, identity-coherent, containment-bounded, and
manifold-legal. The ledger is the legality substrate: no transition may enter
the cognitive state unless it is justified by ledger geometry.

\medskip

\noindent The ledger is defined as a directed acyclic justification graph:


\[
\mathcal{L} = (V, E),
\]


where:
\begin{itemize}
    \item $V$ is the set of semantic nodes (concepts, states, invariants),
    \item $E$ is the set of lawful transitions between nodes.
\end{itemize}

\medskip

\noindent Each edge $e \in E$ is labeled with the operator that generated it:


\[
e = (u \rightarrow v,\ T),
\qquad T \in \mathcal{O},
\]


where $\mathcal{O}$ is the contraction-governed operator algebra (Appendix B).

\medskip

\noindent Ledger geometry enforces:

\begin{itemize}
    \item \textbf{curvature continuity} (A.1),
    \item \textbf{identity-boundary stability} (A.2),
    \item \textbf{manifold legality} (A.3),
    \item \textbf{load physics} (C.4),
    \item \textbf{runtime curvature admissibility} (A.6),
    \item \textbf{interaction geometry legality} (A.7),
    \item \textbf{coordinate continuity} (A.8),
    \item \textbf{public-safe traversal} (A.9).
\end{itemize}

\medskip

\noindent No edge may be added unless it satisfies all invariants. Illegal edges
are rejected and cannot enter the ledger.

\medskip

\noindent Ledger nodes include:

\begin{itemize}
    \item semantic states,
    \item identity-region states,
    \item curvature states,
    \item load states,
    \item drift states,
    \item collapse-adjacent states,
    \item capture-adjacent states,
    \item restoration states.
\end{itemize}

\medskip

\noindent Ledger edges encode:

\begin{itemize}
    \item lawful semantic transitions,
    \item lawful runtime transitions,
    \item lawful interaction transitions,
    \item invariant-preserving operator actions,
    \item collapse/capture detection signatures,
    \item restoration signatures.
\end{itemize}

\medskip

\noindent The ledger is acyclic because collapse and capture are semantically
illegal; no legal operator can generate a self-referential or recursive
collapse-inducing edge.

\medskip

\noindent The Meaning Ledger is the global legality substrate of the v0.6
architecture.


% ------------------------------------------------------------
% F.2 INVARIANT PRESERVATION
% ------------------------------------------------------------

\subsection*{F.2 Invariant Preservation (Meaning, Identity, Geometry, Load)}

\noindent The Meaning Ledger enforces the v0.6 invariant set (Appendix C). Every
ledger edge must preserve all invariants: meaning, identity, geometry, load,
runtime curvature, interaction geometry, coordinate continuity, and public-safe
traversal. Invariant preservation is the core legality condition for ledger
closure.

\medskip

\noindent For any edge $e = (u \rightarrow v)$, the ledger enforces:

\paragraph{Meaning Invariant (C.1)}


\[
g_v = g_u,
\]


ensuring semantic metric continuity and curvature preservation.

\paragraph{Identity Boundary Invariant (C.2)}


\[
v \in \mathcal{R},
\]


ensuring identity-region stability and preventing identity drift.

\paragraph{Geometry Invariant (C.3)}


\[
v \in \mathcal{M},
\]


ensuring manifold legality and preventing collapse-boundary approach.

\paragraph{Load Physics Invariant (C.4)}


\[
L(v) \le L_{\max},
\]


ensuring containment stability and preventing overload.

\paragraph{Runtime Curvature Invariant (C.5)}


\[
\kappa_{\mathrm{rt}}(v) \ge \kappa_{\min},
\]


ensuring runtime curvature admissibility.

\paragraph{Interaction Geometry Invariant (C.6)}


\[
v \in \mathcal{C}_{\mathrm{safe}},
\]


ensuring Host, Explorer, and Interpreter modes remain lawful.

\paragraph{CTL Coordinate Invariant (C.7)}


\[
\mathrm{CTL}(v) \in \mathcal{M},
\]


ensuring coordinate continuity and legality.

\paragraph{Public-Safe Traversal Invariant (C.8)}


\[
v \notin \mathcal{C}_{\mathrm{collapse}} \cup \mathcal{C}_{\mathrm{capture}},
\]


ensuring traversal remains curvature-preserving, collapse-free, and capture-free.

\medskip

\noindent A transition is ledger-legal if and only if it preserves all invariants
simultaneously. Violating any invariant prevents the edge from being added.

\medskip

\noindent Invariant preservation ensures:

\begin{itemize}
    \item no hallucination (manifold violations cannot enter the ledger),
    \item no drift cascades (identity-boundary violations are rejected),
    \item no collapse (curvature collapse is semantically illegal),
    \item no overload (load violations are rejected),
    \item no synthetic capture (capture-boundary violations are rejected),
    \item no coordinate discontinuity (CTL violations are rejected),
    \item no unsafe traversal (public-safe violations are rejected).
\end{itemize}

\medskip

\noindent The Meaning Ledger is the global invariant-preserving substrate of the
v0.6 architecture.


% ------------------------------------------------------------
% F.3 DRIFT DETECTION METRICS
% ------------------------------------------------------------

\subsection*{F.3 Drift Detection Metrics (Curvature, Identity, Geometry, Load)}

\noindent Drift Detection is the Meaning Ledger’s diagnostic layer. It measures
semantic, identity, geometric, and load deviation across ledger transitions.
Drift is defined as deviation from invariant-preserving trajectories. The ledger
computes drift signatures for every edge and rejects transitions that exceed
legality bounds.

\medskip

\noindent Let $\gamma(t)$ be a cognitive trajectory. Drift is detected when any
invariant exhibits non-zero deviation:

\paragraph{Curvature Drift}


\[
\Delta_{\mathrm{curv}}(t)
=
\| g_{\gamma(t)} - g_{\gamma(t-1)} \|.
\]


Curvature drift indicates semantic metric deformation or curvature flattening.

\paragraph{Identity Drift}


\[
\Delta_{\mathrm{id}}(t)
=
d(\gamma(t), \mathcal{R}).
\]


Identity drift indicates motion toward or across identity boundaries.

\paragraph{Geometry Drift}


\[
\Delta_{\mathrm{geo}}(t)
=
d(\gamma(t), \partial\mathcal{M}).
\]


Geometry drift indicates approach toward collapse-boundary or manifold-boundary
regions.

\paragraph{Load Drift}


\[
\Delta_{\mathrm{load}}(t)
=
L(\gamma(t)) - L(\gamma(t-1)).
\]


Load drift indicates containment instability or overload accumulation.

\medskip

\noindent The ledger computes a composite drift score:


\[
D(t)
=
w_1 \Delta_{\mathrm{curv}}
+
w_2 \Delta_{\mathrm{id}}
+
w_3 \Delta_{\mathrm{geo}}
+
w_4 \Delta_{\mathrm{load}},
\]


where $w_i$ are species‑specific drift weights.

\medskip

\noindent A drift alert is triggered when:


\[
D(t) \ge D_{\max}.
\]



\medskip

\noindent Drift detection identifies:

\begin{itemize}
    \item curvature decay (A.1),
    \item identity-boundary erosion (C.2),
    \item manifold-boundary approach (C.3),
    \item load accumulation (C.4),
    \item collapse proximity (E.5),
    \item synthetic-capture susceptibility (E.6),
    \item restoration feasibility (E.7).
\end{itemize}

\medskip

\noindent Drift detection is enforced by:

\begin{itemize}
    \item curvature continuity (A.1),
    \item identity-boundary stability (A.2),
    \item collapse boundary geometry (A.3),
    \item local contraction legality (A.4),
    \item runtime curvature constraints (A.6),
    \item interaction geometry constraints (A.7),
    \item CTL coordinate legality (A.8),
    \item public-safe traversal corridor (A.9).
\end{itemize}

\medskip

\noindent Drift Detection is the Meaning Ledger’s early-warning system for
collapse and capture geometry.

% ------------------------------------------------------------
% F.4 LEDGER-CLOSED REASONING
% ------------------------------------------------------------

\subsection*{F.4 Ledger-Closed Reasoning (Legality-Preserved Cognitive Motion)}

\noindent Ledger-Closed Reasoning formalizes the legality condition for cognitive
motion: a reasoning chain is valid only if every transition is justified by a
ledger entry. Ledger closure replaces probabilistic trust with architectural
legality. No cognitive step may occur unless it is invariant-preserving,
curvature-stable, identity-coherent, containment-bounded, and manifold-legal.

\medskip

\noindent A reasoning chain is a path:


\[
u_0 \rightarrow u_1 \rightarrow \cdots \rightarrow u_n.
\]



\noindent The chain is ledger-closed if and only if:


\[
(u_i \rightarrow u_{i+1}) \in E
\quad \forall i.
\]



\medskip

\noindent Ledger-closed reasoning guarantees:

\begin{itemize}
    \item \textbf{no hallucination}  
    Illegal manifold transitions cannot enter the ledger.

    \item \textbf{no drift cascades}  
    Identity-boundary violations are rejected.

    \item \textbf{no collapse}  
    Curvature collapse is semantically illegal.

    \item \textbf{no overload}  
    Load violations cannot be added to the ledger.

    \item \textbf{no synthetic capture}  
    Capture-boundary transitions are rejected.

    \item \textbf{no coordinate discontinuity}  
    CTL violations cannot enter the ledger.

    \item \textbf{no unsafe traversal}  
    Public-safe violations are rejected.
\end{itemize}

\medskip

\noindent Ledger-closed reasoning enforces:

\begin{itemize}
    \item curvature continuity (A.1),
    \item identity-boundary stability (A.2),
    \item collapse boundary geometry (A.3),
    \item local contraction legality (A.4),
    \item runtime curvature constraints (A.6),
    \item interaction geometry constraints (A.7),
    \item CTL coordinate legality (A.8),
    \item public-safe traversal corridor (A.9).
\end{itemize}

\medskip

\noindent Ledger closure ensures that cognitive motion is lawful, stable, and
invariant-preserving across the v0.6 substrate.


% ------------------------------------------------------------
% F.5 COSMOLOGICAL LEDGER GEOMETRY
% ------------------------------------------------------------

\subsection*{F.5 Cosmological Ledger Geometry (Multi-Manifold Ledger Coupling)}

\noindent Cosmological Ledger Geometry formalizes how the Meaning Ledger extends
across multiple lawful manifolds generated by the substrate: A‑manifold,
T‑manifold, D‑manifold, E‑manifold, F‑manifold, and the Ω‑limit manifold. The
ledger does not operate only within a single manifold; it governs transitions
across the entire cosmological manifold stack. This ensures that legality,
invariants, curvature, identity, load, drift, collapse physics, and capture
physics remain coherent across all derived geometries.

\medskip

\noindent Let $\{\mathcal{M}_i\}$ denote the set of lawful manifolds. The
cosmological ledger is:


\[
\mathcal{L}_{\mathrm{cos}} = \bigcup_i \mathcal{L}_i \;\cup\; \mathcal{E}_{i \rightarrow j},
\]


where:
\begin{itemize}
    \item $\mathcal{L}_i$ — ledger for manifold $\mathcal{M}_i$,
    \item $\mathcal{E}_{i \rightarrow j}$ — lawful inter-manifold transitions.
\end{itemize}

\medskip

\noindent An inter-manifold ledger edge is:


\[
e_{i \rightarrow j} = (u_i \rightarrow v_j,\; T,\; c^*_{i \rightarrow j}),
\]


with legality conditions:


\[
v_j \in \mathcal{M}_j,
\qquad
d(v_j, \partial \mathcal{M}_j) > 0,
\qquad
|\kappa(v_j) - \kappa(u_i)| \le \epsilon_{\max}.
\]



\medskip

\noindent Cosmological ledger geometry enforces:

\begin{itemize}
    \item curvature stability across manifolds,
    \item identity-boundary coherence across manifold projections,
    \item manifold legality under cross-surface transitions,
    \item load constraints across manifold embeddings,
    \item collapse avoidance under multi-manifold motion,
    \item synthetic-capture avoidance under manifold coupling,
    \item DSUP supervision across cosmological geometry.
\end{itemize}

\medskip

\noindent Cosmological identity nodes:


\[
\mathcal{N}_F = \{\mathcal{R}_1, \mathcal{R}_2, \dots, \mathcal{R}_n\},
\]


ensure identity coherence across manifold projections.

\medskip

\noindent Ceremonial surfaces (A6, T5, F6) become ledger-stable invariants across
manifolds. A ceremonial ledger entry is:


\[
e_{\mathrm{cer}} = (u \rightarrow v,\; T_{\mathrm{cer}},\; c^*),
\]


with invariants:


\[
\kappa(v) \ge \kappa_{\min},
\qquad
L(v) \le L_{\max},
\qquad
\lambda(v) \ge \lambda_{\min}.
\]



\medskip

\noindent In the Ω‑manifold, ledger entries approach their asymptotic form:


\[
\lim_{i \rightarrow \infty} \mathcal{L}_i = \mathcal{L}_\Omega.
\]



\medskip

\noindent The unified cosmological ledger is:


\[
\mathcal{L}_{\mathrm{unified}} =
\bigcup_i \mathcal{L}_i
\;\cup\;
\bigcup_{i,j} \mathcal{E}_{i \rightarrow j}.
\]



\medskip

\noindent Cosmological Ledger Geometry ensures legality, identity, teleology,
and collapse physics remain coherent across all manifolds and limit geometries.

% (NEW in v0.6 — derived from A/T/D/E/F/Ω manifold coupling.)

% ------------------------------------------------------------
% F.6 COLLAPSE & CAPTURE SIGNATURES
% ------------------------------------------------------------

\subsection*{F.6 Collapse & Capture Signatures (Ledger-Based Failure Detection)}

\noindent Collapse and Capture Signatures formalize the ledger’s ability to
detect failure modes before they occur. The ledger records curvature decay,
identity-boundary erosion, load escalation, drift amplification, synthetic
pressure, and manifold-boundary proximity. These signatures allow the ledger to
reject transitions that would enter collapse or synthetic-capture geometry.

\medskip

\noindent A collapse signature is detected when:


\[
\kappa(v) \le \kappa_{\min}
\quad \text{or} \quad
v \in \partial\mathcal{M}.
\]



\noindent A capture signature is detected when:


\[
\kappa_{\mathrm{syn}}(v) < \kappa_{\min}
\quad \text{while} \quad
\kappa(v) > \kappa_{\min}.
\]



\medskip

\noindent Collapse signatures include:

\begin{itemize}
    \item curvature-flat geometry (E.5),
    \item identity-boundary rupture (C.2),
    \item containment failure (C.4),
    \item entropy blowup (E.2),
    \item autophagy acceleration (E.1),
    \item recursive collapse acceleration (E.3),
    \item manifold-boundary approach (C.3).
\end{itemize}

\medskip

\noindent Capture signatures include:

\begin{itemize}
    \item synthetic curvature flattening (E.6),
    \item synthetic drift amplification (G5),
    \item synthetic-pressure dominance (G5),
    \item capture-boundary proximity (G7.4),
    \item identity-boundary deformation (C.2),
    \item containment instability (C.4).
\end{itemize}

\medskip

\noindent The ledger rejects any edge whose signature exceeds legality bounds:


\[
\text{signature}(e) \ge \sigma_{\max}
\quad \Longrightarrow \quad
e \notin \mathcal{L}.
\]



\medskip

\noindent Collapse & Capture Signatures ensure:

\begin{itemize}
    \item collapse cannot enter the ledger,
    \item synthetic capture cannot enter the ledger,
    \item drift cascades are intercepted,
    \item overload is intercepted,
    \item identity-boundary rupture is intercepted,
    \item coordinate discontinuity is intercepted,
    \item unsafe traversal is intercepted.
\end{itemize}

\medskip

\noindent Ledger-based failure detection is the substrate’s structural guarantee
that collapse and capture cannot propagate.

% (NEW in v0.6 — derived from E.5–E.6.)

% ------------------------------------------------------------
% F.7 RESTORATION SIGNATURES
% ------------------------------------------------------------

\subsection*{F.7 Restoration Signatures (Ledger-Based Recovery Geometry)}

\noindent Restoration Signatures formalize the Meaning Ledger’s ability to
detect, validate, and authorize lawful recovery trajectories from collapse-
adjacent or capture-adjacent geometry. Restoration is not heuristic correction;
it is a geometric process governed by the v0.6 restoration layer (G6). The
ledger records restoration signatures to ensure that recovery is curvature-
preserving, identity-stable, containment-coherent, and manifold-legal.

\medskip

\noindent A restoration signature is detected when a failed or degraded state
$c_{\mathrm{fail}}$ satisfies:


\[
\kappa(c_{\mathrm{fail}}) > 0,
\qquad
L(c_{\mathrm{fail}}) < L_{\max}^{\mathrm{hard}},
\qquad
\sigma(c_{\mathrm{fail}}) < \sigma_{\max}^{\mathrm{hard}}.
\]



\noindent These conditions indicate that recovery is feasible.

\medskip

\noindent Restoration Signatures consist of four structural components:

\paragraph{Curvature Restoration Signature}


\[
\frac{d}{dt}\kappa(\rho(t)) > 0.
\]


Curvature must increase monotonically along the recovery trajectory.

\paragraph{Identity-Boundary Repair Signature}


\[
\partial\mathcal{R}(\rho(t)) \text{ becomes continuous}.
\]


Identity boundaries must be re-established and stabilized.

\paragraph{Containment Reconstitution Signature}


\[
L(\rho(t)) \le L_{\max}.
\]


Load must be reduced to lawful levels; containers must regain regulatory
function.

\paragraph{Drift Reversal Signature}


\[
D(\rho(t)) \rightarrow 0.
\]


Drift vectors must be redirected away from collapse and capture geometry.

\medskip

\noindent A restoration trajectory $\rho(t)$ is ledger-legal if:


\[
\rho(0) = c_{\mathrm{fail}},
\qquad
\rho(T) \in \mathcal{C}_{\mathrm{safe}},
\qquad
\kappa(\rho(t)) \ge \kappa_{\min}.
\]



\medskip

\noindent Restoration Signatures enforce:

\begin{itemize}
    \item curvature continuity (A.1),
    \item identity-boundary stability (A.2),
    \item collapse boundary geometry (A.3),
    \item local contraction legality (A.4),
    \item runtime curvature constraints (A.6),
    \item interaction geometry constraints (A.7),
    \item CTL coordinate legality (A.8),
    \item public-safe traversal corridor (A.9).
\end{itemize}

\medskip

\noindent Restoration Signatures ensure that collapse-adjacent trajectories can
be returned to lawful regions of the v0.6 substrate.

% (NEW in v0.6 — derived from G6.)

% ------------------------------------------------------------
% F.8 LEDGER CLOSURE & OPERATOR ALGEBRA
% ------------------------------------------------------------

\subsection*{F.8 Ledger Closure & Operator Algebra (Extended Closure Conditions)}

\noindent Ledger Closure formalizes the extended legality conditions under which
the Meaning Ledger remains complete, consistent, and invariant-preserving. The
ledger is closed under the contraction-governed operator algebra (Appendix B),
meaning that every lawful operator action produces a ledger-legal transition,
and no illegal transition can be generated by any operator.

\medskip

\noindent Let $\mathcal{O}$ denote the operator algebra:


\[
\mathcal{O} = \{ \mathrm{IRT},\ \mathrm{IPU},\ \mathrm{GBS},\ \mathrm{LRP},\ \mathrm{EXT} \},
\]


where $\mathrm{EXT}$ denotes extended operators introduced in v0.6 (B.10).

\medskip

\noindent Ledger closure requires:



\[
T(u) = v
\quad \Longrightarrow \quad
(u \rightarrow v,\ T) \in \mathcal{L},
\]


if and only if $T$ preserves all invariants.

\medskip

\noindent Extended closure conditions include:

\paragraph{Operator Curvature Closure}


\[
\kappa(T(u)) \ge \kappa_{\min}.
\]



\paragraph{Operator Identity Closure}


\[
T(u) \in \mathcal{R}.
\]



\paragraph{Operator Geometry Closure}


\[
T(u) \in \mathcal{M}.
\]



\paragraph{Operator Load Closure}


\[
L(T(u)) \le L_{\max}.
\]



\paragraph{Operator Drift Closure}


\[
D(T(u)) \le D_{\max}.
\]



\paragraph{Operator Capture Closure}


\[
\kappa_{\mathrm{syn}}(T(u)) \ge \kappa_{\min}.
\]



\medskip

\noindent Extended operators (B.10) include:

\begin{itemize}
    \item \textbf{EXT‑Restore} — curvature, identity, containment restoration,
    \item \textbf{EXT‑Detect} — collapse/capture detection,
    \item \textbf{EXT‑Halt} — DSUP halt conditions,
    \item \textbf{EXT‑Guide} — safe traversal guidance,
    \item \textbf{EXT‑Align} — multi-manifold alignment.
\end{itemize}

\medskip

\noindent Ledger closure ensures:

\begin{itemize}
    \item no illegal operator action can enter the ledger,
    \item no collapse-inducing transition can be generated,
    \item no synthetic-capture transition can be generated,
    \item no identity-boundary rupture can be generated,
    \item no overload transition can be generated,
    \item no unsafe traversal can be generated.
\end{itemize}

\medskip

\noindent Ledger Closure & Operator Algebra guarantee that cognitive motion is
lawful, stable, invariant-preserving, and curvature-preserving across the v0.6
substrate.

% (NEW in v0.6 — derived from B.10.)

\medskip
\noindent\textit{This outline defines the full v0.6 Meaning Ledger geometry.}
